\chapter{METODOLOGÍA}

\section{Metodología}

La metodología empleada en el desarrollo del proyecto ITPA: Monitorizador IOT se estructuró en base a cuatro aspectos principales: la naturaleza de los datos, la cronología de las actividades, los métodos utilizados y las fuentes bibliográficas. A continuación, se detallan estos elementos.

\subsection{Naturaleza de los Datos}
El proyecto se basó en datos recolectados a través de dispositivos IoT instalados en un entorno controlado (invernaderos). Los datos correspondieron a variables ambientales críticas para el crecimiento vegetal, como temperatura, humedad, niveles de luz y horarios de exposición lumínica. Estos datos se caracterizan por ser:
\begin{itemize}
    \item \textbf{Cuantitativos:} Datos numéricos generados por sensores, como lecturas de temperatura (\textdegree C) y humedad (\%).
    \item \textbf{Continuos:} Capturados en tiempo real, permitiendo análisis históricos y patrones temporales.
    \item \textbf{Multidimensionales:} Incluyen múltiples variables relacionadas con las condiciones ambientales.
\end{itemize}

\subsection{Cronología}
El desarrollo del proyecto se llevó a cabo en el periodo comprendido entre septiembre y noviembre, dividido en tres fases principales:
\begin{itemize}
    \item \textbf{Fase 1 (Septiembre):} Investigación bibliográfica, definición de objetivos y selección de tecnologías base.
    \item \textbf{Fase 2 (Octubre):} Desarrollo del framework IoT, integración de sensores y pruebas preliminares en el sistema de invernaderos.
    \item \textbf{Fase 3 (Noviembre):} Validación del sistema, análisis de datos recolectados y creación de modelos gráficos descargables.
\end{itemize}

\subsection{Métodos Empleados}
Se utilizaron metodologías ágiles para gestionar el desarrollo del proyecto, particularmente \textbf{Scrum}, lo que permitió iteraciones rápidas y la incorporación de retroalimentación constante. Las actividades principales incluyeron:
\begin{itemize}
    \item \textbf{Desarrollo del Framework:} Se emplearon Node.js, TypeScript y MongoDB para la construcción de una arquitectura modular y escalable.
    \item \textbf{Recolección de Datos:} Uso de sensores IoT configurados con dispositivos Raspberry Pi para capturar y transmitir datos en tiempo real.
    \item \textbf{Análisis de Datos:} Implementación de algoritmos para agrupar por fechas, crear modelos predictivos y graficar las variables ambientales.
    \item \textbf{Validación:} Pruebas funcionales y de estrés para garantizar la robustez y fiabilidad del sistema.
\end{itemize}

\subsection{Fuentes Utilizadas}
El proyecto se apoyó en referencias bibliográficas relevantes que aportaron conocimientos clave sobre la integración de IoT y el uso de iluminación artificial en la agricultura. Entre las principales fuentes destacan:
\begin{itemize}
    \item Fernández, J. (2021). \textit{Scalable IoT Solutions for Agriculture and Industry}. TechPress.
    \item González, M., \& Hernández, J. (2019). \textit{Artificial Lighting for Enhanced Plant Growth}. \textit{Plant Science Journal}, 40(3), 200-210.
    \item López, E., \& Torres, M. (2020). \textit{IoT Frameworks for Data Monitoring in Agriculture}. \textit{International Journal of IoT Systems}, 12(2), 150-160.
    \item Zhang, W., \& Liu, H. (2021). \textit{LED Lighting in Modern Agriculture: Applications and Benefits}. \textit{Journal of Agricultural Engineering}, 45(4), 300-310.
\end{itemize}

Esta metodología permitió desarrollar un sistema eficiente y adaptado a las necesidades del proyecto, asegurando su éxito y aplicabilidad en el ámbito agrícola.
